\documentclass[11pt,a4paper]{extarticle}

\usepackage[utf8]{inputenc}
\usepackage[T1]{fontenc}
\usepackage[brazil]{babel}
\usepackage{amsmath, amssymb, amsthm}
\usepackage[margin=1.5cm, top=2.6cm, headheight=30pt]{geometry}
\usepackage{xcolor}
\usepackage{microtype}
\usepackage{enumitem}
\usepackage{fancyhdr}
\usepackage{tcolorbox}
\usepackage{pgffor}

\definecolor{azulescuro}{RGB}{20, 50, 100}
\definecolor{laranja}{RGB}{230, 120, 30}

% Valores das questões (ajuste aqui)
\newcommand{\valorA}{3,5}
\newcommand{\valorB}{3,5}
\newcommand{\valorC}{3,0}

% Linhas para resposta
\newcommand{\linhas}[1]{%
  \par\nointerlineskip
  \foreach \i in {1,...,#1}{%
    \vspace{0.6cm}\noindent\textcolor{black}{\rule{\linewidth}{0.5pt}}\par\nointerlineskip%
  }%
}

% Cabeçalho de questão
\newcommand{\questao}[2]{%
  \vspace{0.6em}
  \begin{tcolorbox}[colback=azulescuro!6, colframe=azulescuro, boxrule=0.6pt,
    arc=1.5pt, left=6pt, right=6pt, top=3pt, bottom=3pt]
    \textbf{\textcolor{azulescuro}{Questão #1}} \hfill \textcolor{laranja}{\textbf{(valor: #2)}}
  \end{tcolorbox}
  \vspace{0.2em}
}

\setlist{nosep}
\setlength{\parindent}{0pt}
\setlength{\parskip}{0.5em}

\pagestyle{fancy}
\fancyhf{}
\renewcommand{\headrulewidth}{0.6pt}
\fancyhead[L]{\small\textbf{BCC502} -- Metodologia Científica para Ciência da Computação}
\fancyhead[R]{\small DECOM/ICEB/UFOP}
\fancyfoot[C]{\small Página \thepage\ de \pageref{LastPage}}
\usepackage{lastpage}

\begin{document}

\begin{center}
  {\Large\bfseries\textcolor{azulescuro}{Extra 1}}\\[0.3em]
  {\large \textbf{BCC502} -- Metodologia Científica para Ciência da Computação}
\end{center}

\vspace{0.3em}
\begin{tcolorbox}[colback=white, colframe=azulescuro!60, boxrule=0.5pt, arc=1.5pt]
  \textbf{Nome:} \rule{8.2cm}{0.3pt} \hfill \textbf{Matrícula:} \rule{3cm}{0.3pt}\\[0.9em]
  \textbf{Data:} \rule{1.6cm}{0.3pt}/\rule{1.6cm}{0.3pt}/\rule{2.2cm}{0.3pt}
  \hfill \textbf{Nota:} \rule{2.5cm}{0.3pt}
\end{tcolorbox}

\textbf{Instruções.} Responda a cada questão no espaço reservado, com letra legível e por extenso. Respostas que apenas listem termos, sem articular o raciocínio, receberão nota parcial. Não é permitida a consulta a qualquer material.

% ------------------------------------------------------------
1. Quais problemas do indutivismo o falsificacionismo resolve? Como ele resolve? O que ainda fica em aberto?
\linhas{23}
\newpage
% ------------------------------------------------------------
2. Por que o falsificacionismo não explica a ciência e como Lakatos e Kuhn tentam fechar essa lacuna?
\linhas{15}

% ------------------------------------------------------------
3. Para que serve o conhecimento sobre filosofia da ciência para não cientistas?

\linhas{15}

\end{document}